% ----------------------------------------------------------
\chapter{Conclusão}
% ----------------------------------------------------------

O presente trabalho teve como objetivo o desenvolvimento de uma plataforma computacional para o pré-dimensionamento e otimização de vigas de madeira, integrando modelagem paramétrica, verificação normativa e técnicas de otimização multiobjetivo. Os resultados obtidos demonstram que os objetivos propostos foram plenamente atingidos, tanto do ponto de vista computacional quanto do ponto de vista estrutural.

No âmbito do projeto paramétrico, a ferramenta desenvolvida mostrou-se eficiente e funcional, permitindo ao usuário a obtenção de soluções de projeto em um curto intervalo de tempo, por meio de uma interface intuitiva e de fácil operação. Essa característica representa um avanço significativo em relação aos métodos tradicionais de dimensionamento, ao possibilitar a geração rápida de múltiplas alternativas de projeto, favorecendo análises comparativas e decisões técnicas mais fundamentadas.

Em termos de otimização, o algoritmo NSGA-II demonstrou elevada eficiência na obtenção da fronteira de Pareto, identificando soluções estruturalmente viáveis e materialmente eficientes. A comparação direta com o método de Monte Carlo evidenciou de forma clara a superioridade do processo de otimização, uma vez que as soluções obtidas por amostragem aleatória apresentaram comportamento mais conservador, com maior consumo de material e baixa utilização da capacidade resistente da estrutura.

A análise mostrou que o modelo responde de forma coerente às diferentes cargas de projeto, enquanto o tabuleiro da classe TB-450 demandou seções mais robustas e maior volume de madeira para garantir baixa deformação sob altas exigências, o tabuleiro TB-240 apresentou um dimensionamento mais esguio e econômico, operando próximo aos limites normativos de flecha. Essa proporcionalidade entre o aumento das solicitações e a resposta estrutural valida a precisão física do sistema desenvolvido.

Além disso, vale mencionar que o dimensionamento eficiente de longarinas e o tabuleiro não pode ser tratado de forma isolada no contexto de estruturas reais, como pontes de madeira. Elementos como as ligações estruturais, os sistemas de fundação e os pilares exercem influência direta no comportamento global da estrutura, sendo fundamentais para a segurança, durabilidade e desempenho estrutural do sistema como um todo. Assim, a integração futura desses componentes ao modelo computacional representa uma perspectiva natural de evolução da ferramenta desenvolvida.

Por fim, conclui-se que a plataforma proposta constitui uma ferramenta robusta e tecnicamente consistente para apoio ao projeto estrutural em madeira, contribuindo para a modernização dos processos de dimensionamento, para a racionalização do uso de material e para a difusão de metodologias baseadas em otimização multiobjetivo no contexto da engenharia estrutural. O trabalho demonstra o potencial da combinação entre modelagem paramétrica, computação científica e técnicas de otimização como suporte à tomada de decisão em projetos estruturais, abrindo caminho para aplicações mais amplas e integradas em sistemas estruturais complexos.

